
The standard is not about event logs and not about configuration databases,
and the quickest way to show that is to run it on something else. UCI Adult
(Census Income) predicts whether a person's income exceeds a threshold from
\exRows\ records. Suppose an analyst wants to report what \texttt{occupation}
is worth. It is the expensive field: a survey has to ask about it and code it,
where age, sex, race and country come with the sampling frame and education
comes with an enrolment record.

Worth compared with what? Against the frame alone, against the frame plus
education, or against the frame plus education plus hours worked. Measured
how, at which operating point, and with the register how complete. Over the
resulting \exCells\ cells the reduction runs \exRmin\ to \exRmax. The baseline
axis and the metric axis carry the same first-order share of the variance in
the increment to three figures --- \exSbaselinePct\ and \exSmetricPct\ ---
against \exSpopulationPct\ for the register population, and a one-number
report misstates the sign for a mean \exMisreportPct\ of the other cells. That
the two lead axes tie is a property of this dataset and not a general one; it
is the sort of thing a surface reports and a single number cannot.

\texttt{examples/worked\_example.py} produces all of it in \exampleLines\
statements against the package's public interface, and ends by calling
\texttt{float()} on the surface so a reader can see what the refusal looks
like.

